\scanpage{172}
C'est dans ce dernier diagramme qu'interviennent les ``nouveaux'' topos élémentaires qu'on vient d'introduire, sous la forme des $\unsure{\mathcal V}^{\Delta_1}_{\mathbb P,X}$ et $\unsure{\mathcal V}^{\Delta_2}_{\mathbb P,X}$.

Je voudrais maintenant expliciter la variance des topos $\unsure{\mathcal V}^{d'_0}_{d_*,d''_0}$ pour un \deleted{triplet $(d''_*\subset d'_*\subset d_*)$}\inserted{$d'_0,d''_0\in d_*,\ d'_0\geq d''_0$} variable, et celle des $\Pi_1$ associés, et expliciter comment, à l'aide de ces topos ou groupoïdes, on peut reconstituer tous les ``topos canoniques'' associés à la stratification envisagée -- ou leurs groupoïdes fondamentaux.

Avant \deleted{de \unsure{voir}} d'examiner cependant la variance des $\unsure{\mathcal V}^{d'_0}_{d_*,d''_0}$, je voudrais \unsure{regarder} un peu la structure de ces ``\unsure{espaces|espèces}''.

\unsure{Heuristique|Heuristiques}, prenons le fibré normal de $D_{d_*}$ \deleted{dans $D_{d''_0}$, soit $\check N_{d_*,\unsure{d''_*|d''_0}}$}
\deleted{\[
(282)\qquad\check N_{d_*,d''_0}\overset{\mathrm{def}}=\check N_{D_{d_*},D_{d''_0}}\qquad\uline{\text{fibré normal}}
\]}
\uline{NB}
\margin{171}
